\source{Лекции, с. 28–29; строгие условия связности и детали индукции восстановлены составителем; Карпов, §1.9}
\begin{definition}[Эйлеров путь и цикл]
Эйлеров путь проходит по каждому ребру графа ровно один раз.
Если он замкнут, это эйлеров цикл; граф с таким циклом называется эйлеровым.
Вершины могут повторяться. В орграфе дуги проходят только по направлению.
Далее существенная часть — вершины ненулевой степени со всеми рёбрами.
В лекциях исходный граф предполагается связным; отделение существенной
части позволяет также учитывать изолированные вершины.
\end{definition}
\begin{theorem}[Критерии Эйлера]
Для графа с хотя бы одним ребром необходима связность существенной части
после забывания направлений. При этом условии верны следующие критерии:
\begin{center}
\begin{tabular}{p{.13\linewidth}|p{.34\linewidth}|p{.39\linewidth}}
 & Замкнутый обход & Открытый обход с концами \(s\ne t\)\\\hline
Граф & Все степени чётны & Ровно две нечётные вершины \(s,t\)\\[3pt]
Орграф & \(d^+(v)=d^-(v)\) для всех \(v\) &
\(d^-(s)-d^+(s)=1\), \(d^+(t)-d^-(t)=1\); у остальных баланс нулевой
\end{tabular}
\end{center}
\end{theorem}
\begin{proof}
Необходимость: у внутренних посещений вход и выход образуют пару.
У начала открытого пути остаётся один лишний выход, у конца — вход.
Один обход всех рёбер также даёт требуемую связность.

\textit{Неориентированный цикл: лекционная индукция по числу рёбер.}
В графе с положительными чётными степенями есть простой цикл:
продолжаем простой путь до первой вынужденной встречи с прежней вершиной.
Удаляем рёбра этого цикла \(C\). В каждой оставшейся существенной
компоненте все степени чётны и рёбер меньше; по индукции есть эйлеров
цикл. Каждая такая компонента имеет общую вершину с \(C\): иначе исходная
связность нарушалась бы (все удалённые рёбра имеют концы на \(C\)).
Начинаем обход \(C\), а при посещении общей вершины вставляем обход
соответствующей компоненты и возвращаемся к обходу \(C\).
Получено прохождение всех рёбер ровно один раз. Если остаточных рёбер нет,
сам \(C\) завершает доказательство; это также база индукции.

\textit{Ориентированный цикл.} Пока есть дуги, в сбалансированном орграфе
есть направленный цикл. Действительно, после входа в вершину у неё есть
исходящая дуга, а бесконечный выбор продолжений в конечном множестве
даёт цикл. Удаление дуг простого направленного цикла сохраняет баланс.
Остаточные компоненты графа с забытой ориентацией сбалансированы;
по индукции у них есть направленные эйлеровы циклы. Склейка с исходным
циклом проводится так же. Здесь сильная связность не предполагается:
она сама следует из полученного обхода всех ненулевых вершин.

\textit{Открытый обход.} В неориентированном случае добавим ребро \(st\),
в ориентированном — дугу \(t\to s\). Получится граф с условием для цикла.
Строим этот цикл, затем разрезаем его запись по добавленному ребру/дуге.
Остаётся эйлеров путь из \(s\) в \(t\). Временные кратные рёбра допустимы:
доказательство индукцией различает рёбра, а не только их концы.
\end{proof}
\begin{cor}[Разбиение на пути]
Если связный неориентированный граф имеет \(2k>0\) нечётных вершин,
его рёбра разбиваются на \(k\) непересекающихся по рёбрам открытых путей.
Меньшим числом путей обойти все рёбра невозможно.
\end{cor}
\begin{proof}
Попарно соединим нечётные вершины \(k\) временными рёбрами.
Полученный связный мультиграф имеет все степени чётными. Разрезая его
эйлеров цикл по временным рёбрам, получаем \(k\) путей исходного графа.
Временные рёбра не могут следовать подряд, поскольку имеют разные концы;
поэтому все полученные пути непустые. Каждая нечётная вершина должна
быть концом нечётного числа путей разбиения. У \(r\) путей только \(2r\)
концов, следовательно, \(r\ge k\).
\end{proof}
Если все степени чётны и есть ребро, достаточно одного замкнутого пути.
Для одноточечного графа без рёбер нулевой путь считается эйлеровым.
Петли, если они специально разрешены, дают две инцидентности в неориентированном
случае и по одному входу и выходу в ориентированном; критерии сохраняются.
